\subsection{One-space axioms and regular partitions}

\begin{saresult}{A.301}{Refuted}{The literal simultaneous compact axioms do not characterize o-minimal structures}
If the symbol $R$ denotes an ordered real closed field with its order topology,
the simultaneous requirements M0 and M9 of cote 120 are inconsistent.
Consequently M0--M11 with unary finiteness cannot, as one simultaneous axiom
system, characterize compact definable sets of structures with o-minimal open core.
\end{saresult}
\begin{proof}
M0, at 120/007, is an optional axiom requiring $R\in M_1$ when adopted.
M9, at 120/017, requires all
members of $M_n$ to be compact. The line $R$ is not compact: the open
intervals $(-a,a)$, $a>0$, cover it and no finite subfamily does. If
``compact'' is replaced by definably compact over a non-Archimedean
field, the contradiction persists because $R$ is unbounded. The family
of compact definable sets itself fails M0. These are different drafts:
120/021 explicitly discards the earlier topological structures before
introducing the field operations. Keeping that distinction is necessary.
\end{proof}

\begin{saresult}{A.302}{Proved}{An exact completion criterion for a one-space family}
Let $K$ be a real closed field, $X$ a set, $\mathcal F\subseteq\mathcal P(X)$,
and $j:X\hookrightarrow K^m$ an injection. Form the smallest all-arity
family $\mathcal D=(\mathcal D_n)$ containing all semialgebraic sets over
$K$ and the predicates $j(A)$ for every $A\in\mathcal F$, and closed under
finite Boolean operations, products, coordinate permutations and coordinate
projections. There is an o-minimal expansion of the ordered field $K$ in
which $j(X)$ is definable and its definable subsets are exactly $j(\mathcal F)$
if and only if
\[
\begin{gathered}
 X\in\mathcal F,\qquad
 \mathcal D_1=\{\text{finite unions of points and intervals in }K\},
 \\
 \{D\cap j(X):D\in\mathcal D_m\}=j(\mathcal F).
\end{gathered}
\]
Here parameters from $K$ are allowed. If one-space topology is part of
the data, $j$ must in addition be a homeomorphism onto its image.
The criterion applies whether or not the one-space finite-dimension
axiom $(e_n)$ is initially imposed. It neither derives the missing
coordinates from $(e_n)$ nor substitutes that bound for the all-arity
and trace tests.
\end{saresult}
\begin{proof}
If the expansion exists, its definable sets contain every generator and
are closed under all the listed operations. Thus they contain every
$\mathcal D_n$. Unary finiteness follows, and the trace condition follows
in one direction from this containment and in the other from the named
predicates. Conversely expand the ordered field by one relation symbol
for every $j(A)$. Every formula in this language uses finitely many
predicates. Induction on formulas identifies its definable sets with
$\mathcal D$: conjunction and negation are Boolean operations, adding
unused variables is a product, and existential quantification is
projection; coordinate identifications are intersections with
semialgebraic diagonals followed by projections. Thus the unary condition
is precisely o-minimality, and the trace condition gives exactly the
prescribed family on $X$. The topological assertion is then the specified
compatibility of $j$.
\end{proof}

This criterion exposes the missing data in 121/004--024: one needs
coordinates and all finite-arity existential consequences, including a
test that they create no new subsets of $X$. Conditions imposed on a
Boolean algebra on one space alone cannot supply them. For example,
$\mathcal F=\{\varnothing,X\}$ on an infinite connected compact manifold
satisfies the elementary Boolean, closure, frontier and connected-component
conditions on this family. It cannot be the family of all parameter-definable
subsets in any injective field chart, because those include every singleton.
We do not conflate the full all-arity completion with the open core.

\begin{saresult}{A.309}{Proved}{Dimension bounds and the connectedness convention}
A definable realization of dimension $d$ satisfies the one-space
axiom $(e_d)$ when ``rare'' means nowhere dense. Over a real closed
field which is not Dedekind complete, ordinary connected-component
finiteness instead of definable connected-component finiteness
forces every definable realization in a finite Cartesian power
of the field to be finite.
\end{saresult}
\begin{proof}
For a realized definable set of dimension $d$, the dimension-drop
property gives $(e_d)$ in the definable interpretation: in a chain
of relatively closed definable sets, each nowhere-dense inclusion
lowers dimension, so a chain of $d+1$ such inclusions is empty.
The same assertion holds for locally definable germs on arbitrary
opens, since a finite chain can be represented in one common
definable neighbourhood of each point. Here ``rare'' means
nowhere dense, as in 121/024. If ``non dense'' on 121/006 were read
as merely not dense, the sequence $[0,1/j]$ would refute the
axiom even in the real semialgebraic line.

Condition f) must also be interpreted explicitly over $K$.
In the definable category connected components mean definably
connected components. If f) instead means ordinary topological
components and $K$ is not Dedekind complete, a realization
satisfying it must be finite. Indeed choose an unrealized bounded
cut of $K$. Translating and scaling a bounded interval containing
that cut places an unrealized cut strictly between any two
specified elements of $K$. Its lower and upper parts are disjoint
open subsets covering $K$, and separate these two elements.
Every connected subset of $K$ is therefore a singleton; coordinate
projection gives the same assertion for every subset of $K^m$.
Every positive-dimensional definable set contains a cell with
infinitely many points, so its ordinary components cannot be
finite in number. For a finite set, all subsets and all component
conditions hold. The only Dedekind-complete ordered field is
$\R$: completeness gives the Archimedean property by taking a
supremum of the integers, and the rational cuts then give its
unique ordered-field identification with $\R$. This explains
precisely why ordinary and definable connectedness agree in the
real setting but cannot silently be exchanged over arbitrary
real closed fields.
\end{proof}

\subsection{All closed saturated unions, rather than only closures of strata}

We isolate the purely finite argument needed at 121/018. A finite
semialgebraic partition is \emph{cellular} here if its strata are locally
closed and the closure of every stratum is a union of strata. It is
$F$-regular if for every closed union $Y$ of strata, the nonsmooth locus
$\operatorname{Sing}(Y)$ is also a union of strata. Smooth means locally
a $C^\infty$ embedded submanifold, with dimension locally allowed to vary.
This is the condition on all saturated closed unions asked for in A.3(d).

\begin{saresult}{A.303}{Proved}{Existence and uniqueness of the coarsest $F$-regular partition}
Let $X\subset\R^N$ be semialgebraic and let $\Phi$ be a finite family of
semialgebraic subsets. Assume the standard semialgebraic facts that finite
families admit finite smooth frontier partitions and that the nonsmooth
locus of a closed semialgebraic set is a closed semialgebraic subset of
smaller dimension. Then there is a unique coarsest finite cellular
$F$-regular partition of $X$ adapted to $\Phi$. Its strata are smooth.
No connectedness or equisingularity assertion is included in this result.
\end{saresult}
\begin{proof}
We first produce one such partition, taking care of every closed union.
Start with a smooth frontier partition adapted to $\Phi$. We freeze its
strata successively in decreasing dimension. Suppose strata of dimension
greater than $k$ have already been frozen, and a frontier partition of
the remaining set has been chosen compatibly with their closures. For
each union $A_J$ of closures of frozen strata, and each current
$k$-dimensional stratum $S$, partition $S$ according to the finitely many
sets $\operatorname{Sing}(A_J)\cap S$. Remove from $S$ their relative
boundaries. Each removed boundary has dimension less than $k$, and each
remaining sign-pattern piece is open in $S$ and smooth. We maintain the
invariant that the unfrozen remainder is closed in $X$. At this stage
every current $k$-stratum is open in that remainder, by the frontier
condition and maximality of $k$. Hence the newly frozen pieces are open
in the remainder and their complement is still closed in $X$. Replace
that lower-dimensional complement by a smooth frontier partition adapted
to all frozen closures and to the previous partition. Its strata cannot
accumulate onto a frozen stratum because the remainder is closed and
disjoint from the frozen strata. There are finitely many
steps and finitely many pieces at each step.

Let $S$ be a final $k$-stratum and $Y$ a closed union of final strata.
If $S\not\subset Y$, it is disjoint from $Y$. Otherwise, near a point
$s\in S$, the only other strata of $Y$ that can occur have dimension
greater than $k$ and have closures containing all of $S$. Indeed the
frontier condition forces a closure meeting $S$ to contain $S$, and a
proper frontier has strictly smaller dimension. The finitely many
closures not meeting $S$ can be avoided in a neighborhood of $s$.
If no higher stratum of $Y$ approaches $S$, the germ of $Y$ is the smooth
germ of $S$. Otherwise that germ is the germ of the union $A_J$ of the
closures of the approaching higher strata of $Y$. These strata were
already frozen when $S$ was constructed. Consequently the construction
made $\operatorname{Sing}(A_J)\cap S$ either all of $S$ or empty.
The same holds for $\operatorname{Sing}(Y)\cap S$. This proves
$F$-regularity for every closed saturated union simultaneously.

For coarseness let $\mathcal B_P$ denote the finite Boolean algebra of
unions of strata of a cellular $F$-regular partition $P$. It is closed
under taking closure in $X$ and under $Y\mapsto\operatorname{Sing}(Y)$
for closed $Y$ in the algebra. For two such partitions $P,Q$ adapted to
$\Phi$, the algebra $\mathcal B_P\cap\mathcal B_Q$ has the same
properties and contains $\Phi$. Its atoms form a finite semialgebraic
partition $R$. We check cellularity rather than assuming it.

Every atom $A$ has closure in the algebra. It is open in that closure.
To see this, its frontier $\overline A\setminus A$ is in the algebra;
the closure of that frontier is in the algebra too. If it met $A$, it
would contain all of $A$, by atomicity. This is impossible because
$\dim(\overline{\overline A\setminus A})<\dim A$ for a semialgebraic
set. Thus $A$ is disjoint from that closed frontier closure, proving
local closedness. The closure of each atom is a union of atoms, so $R$
is cellular. The singular locus of every closed union is in the algebra,
so $R$ is $F$-regular. In particular, $\operatorname{Sing}(\overline A)$
either contains $A$ or misses it. It cannot contain $A$, which is dense
in its closure, since the singular locus is closed with empty interior.
Thus $A$ is smooth, being open in the smooth locus of $\overline A$.

Among the nonempty class of adapted $F$-regular partitions choose $P$
with the fewest strata. The common coarsening just constructed for $P$
and any other $Q$ cannot have fewer atoms. Hence it equals $P$, and $P$
is coarser than $Q$. This proves existence of a coarsest partition and
its uniqueness.
\end{proof}

The smooth-locus input is recorded in \cite[\S2.A, pp.~635--636]{Fernando2018};
it concerns the smooth locus, not the potentially smaller algebraic regular
locus. Finite smooth frontier partitions and frontier dimension are
\cite[Theorems 1.3 and 1.5, the paragraph following Theorem 1.3,
and Proposition 1.9, pp.~3--5]{Coste2007}. Thus the explicitly stated
standard inputs hold in the semialgebraic setting of this theorem. In particular,
checking only $\operatorname{Sing}(\overline S)$ for individual strata,
as the short argument at 121/018 does, is not the proof supplied here.

\begin{saresult}{A.304}{Proved}{Connected and pure-dimensional versions of the coarsest partition}
In A.303 one may require all strata to be connected, or all strata
to have a single dimension, or both. Each of the resulting classes
has its own unique coarsest adapted cellular $F$-regular partition.
The connected version also has pure-dimensional strata.
\end{saresult}
\begin{proof}
Existence in the connected class follows by taking connected components
of the finite pieces at each stage of the construction in A.303.
They remain finite and semialgebraic, and the same lower-dimensional
refinement makes their closures saturated. The construction already
gives pure-dimensional pieces.

We check that the common-coarsening argument preserves the extra
requirements. If every stratum of $P$ is connected, every connected
component of a union of its strata is itself a union of strata:
a connected stratum cannot meet two different components. Thus
$\mathcal B_P$ is closed under taking individual connected components.
The intersection $\mathcal B_P\cap\mathcal B_Q$ has this property too,
so its atoms are connected. A connected smooth manifold has constant
dimension, proving the last assertion.

For the pure-dimensional class without connectedness, use for every
closed saturated $Y$ the closed semialgebraic subset
$D_k(Y)=\{y\in Y:\dim_yY\ge k\}$. It is saturated for any cellular
partition into pure-dimensional strata: at a point of a stratum $S$,
the local dimension of $Y$ is the largest dimension of a stratum $T$
of $Y$ for which $S\subseteq\overline T$. The frontier condition
makes this number independent of the point of $S$. Hence both Boolean
algebras, and their intersection, are closed under $D_k$ on closed
members. For an atom $A$ of the intersection, the sets
$A\cap D_k(\overline A)$ are therefore either empty or all of $A$.
Since $A$ is open in its smooth closure locus by A.303, its manifold
dimension is constant. The minimum-number-of-atoms argument of A.303
now applies in each of the indicated classes and proves uniqueness.
\end{proof}

\subsection{The all-arity repair and the compact open core}

\begin{saresult}{A.305}{Proved}{The coherent all-arity algebraic system gives full o-minimality}
Let $K$ be a real closed field. For each $n\geq0$ let
$\mathcal M_n\subseteq\mathcal P(K^n)$ contain $K^n$, the empty set,
and every parameter point. Assume closure under finite intersections,
products, coordinate permutations and coordinate images, and that the
graphs of addition and multiplication belong to the appropriate families.
Call a map moderate when its graph belongs to the corresponding family.
Assume, as in the algebraic zero-set axiom of cote 120, that for every
$Y\in\mathcal M_n$ there are finitely many moderate functions
$f_i:K^n\to K$ with $Y=\bigcap_i f_i^{-1}(0)$.
Then the $\mathcal M_n$ are exactly the sets definable, with parameters,
in an expansion of the ordered field $K$. This expansion is o-minimal
if and only if $\mathcal M_1$ consists of finite unions of points and
intervals. No compactness axiom is imposed in this assertion.
Conversely, the definable sets of every expansion of $K$, with all
definable maps allowed, satisfy the stated zero-set axiom.
\end{saresult}
\begin{proof}
The graph operations in the hypotheses give graphs of polynomial
expressions and inverse images under moderate maps: intersect the
relevant graphs and project away their intermediate values. If
$Y=\{f_1=\cdots=f_s=0\}$, then
\[
 K^n\setminus Y
 =\operatorname{pr}_{K^n}
 \left\{(x,t_1,\ldots,t_s):\sum_{i=1}^s t_i f_i(x)=1\right\}.
\]
Indeed a nonzero value of one $f_i$ supplies the corresponding inverse
$t_i$, and if all values vanish the displayed equation has no solution.
Thus complements belong to the families. The empty list of equations
gives $Y=K^n$ and the complement is empty, as required. Unions now follow
by De Morgan's law. The order is definable because
$x\leq y$ if and only if $y-x=z^2$ for some $z\in K$.
Expand the field by all the predicates in the families. Induction on
first-order formulas, exactly as in A.302, shows that no new sets are
created: negation is complement and existential quantification is a
coordinate image. The converse containment holds because every original
set has been named. The unary condition is the definition of
o-minimality. Conversely, for a definable $Y\subseteq K^n$ the function
which is $0$ on $Y$ and $1$ on its complement is definable and has zero
set $Y$. This last step uses all definable maps, not only continuous
ones. These are the two topological and algebraic drafts distinguished
at 120/021.
\end{proof}

\begin{saresult}{A.306}{Proved}{What compact predicates determine, and an exact realizability criterion}
For an arbitrary expansion $\mathcal R$ of an ordered real closed field
$K$, let $\mathcal C_n$ be all its definable closed bounded subsets of
$K^n$, allowing parameters. The expansion of the ordered field by all
the predicates in all $\mathcal C_n$ is interdefinable with the open
core of $\mathcal R$. Here the open core means the expansion of the
ordered field by every open $\mathcal R$-definable subset of every
$K^n$. Over $\R$, closed bounded has its ordinary compact meaning.

Consequently, a prescribed family $(\mathcal C_n)_n$ of closed bounded
sets is exactly the family of compact sets of an expansion with
o-minimal open core if and only if the minimal all-arity definable
expansion $K_{\mathcal C}$ generated by those predicates is o-minimal
and its closed bounded definable sets are exactly the prescribed
$\mathcal C_n$. This criterion includes the requirement that all
parameter singletons occur.
\end{saresult}
\begin{proof}
Every closed bounded predicate is definable from its open complement,
so it is definable in the open core. For the converse use the
semialgebraic homeomorphism
\[
 \rho:K\longrightarrow(-1,1),\qquad
 \rho(x)=\frac{x}{1+|x|},
 \quad\rho^{-1}(u)=\frac{u}{1-|u|}.
\]
If $U\subseteq K^n$ is open and definable, then $\rho^n(U)$ is open
in $K^n$ and contained in $(-1,1)^n$. Hence
$C=[-1,1]^n\setminus\rho^n(U)$ is closed, bounded and definable.
The equation $U=(\rho^n)^{-1}([-1,1]^n\setminus C)$ recovers $U$
from this one compact predicate and the ordered-field operations.
The case $n=0$ consists only of the empty set and the point and obeys
the same conclusion.

If the requested expansion $\mathcal R$ exists, the first paragraph
identifies $K_{\mathcal C}$ with its open core. It is therefore
o-minimal. Its closed bounded definable sets are exactly $\mathcal C$:
they are definable in $\mathcal R$, and every member of $\mathcal C$
was named. Conversely take $\mathcal R=K_{\mathcal C}$. The stated
conditions give the required compact sets and an o-minimal expansion;
its open core is itself. To verify the last assertion, in an o-minimal
expansion every definable set is a finite union of locally closed cells,
and each locally closed set is an intersection of a closed and an open
definable set, hence definable in the open core. Equivalently the first
paragraph applies to its compact predicates.
\end{proof}

The distinction is strict even over the real numbers. The dense pair
$(\R,\R_{\mathrm{alg}})$ has exactly the semialgebraic open sets by
\cite[Theorem 5 and its proof in \S4.8, pp.~63 and 77;
Corollary 4.9, pp.~77--78]{VandenDriesDensePairs1998}.
Its compact definable sets are therefore exactly the compact
semialgebraic sets: take complements in one direction and use the
field reduct in the other. Yet its named proper dense subfield
$\R_{\mathrm{alg}}$ is not a finite union of points and intervals, so
the dense pair is not o-minimal. Thus compact data cannot distinguish
these two expansions. This is an application of the opened theorem,
not an assumption that unary tameness of open sets implies full
o-minimality of the ambient expansion.

\begin{saresult}{A.307}{Refuted}{An arbitrary prescribed triangulation need not extend}
The last sentence of 121/065, read literally as extension of a given
triangulation, is false already for compact semialgebraic pairs. It
remains false when linear subdivisions of the prescribed triangulation
are allowed. Therefore the compact systems of 121/054--073, if they
include that sentence, cannot characterize all compact definable sets
of any o-minimal expansion of a real closed field.

The true replacement is existence of a triangulation compatible with
any finite family of definable subsets. A prescribed triangulation of
a closed subset can be extended, up to subdivision, if its transition
to the restriction of an ambient triangulation is PL.
\end{saresult}
\begin{proof}
We work over $\R$; this one example suffices to refute the universal
claim. Put
\[
 \begin{split}
 Y&=[-2,2]^2\times\{0\},\\
 X&=Y\ \cup\
 \{(x,y,z):x^2+y^2=1,\ 0\leq z\leq1\}.
 \end{split}
\]
These are compact semialgebraic sets and $Y$ is closed in $X$.
Give $Y$ any ordinary finite straight triangulation and prescribe its
identity realization. Let
$C=\{(x,y,0):x^2+y^2=1\}\subset\operatorname{int}Y$.
At each point of $C$, the germ of $X$ is a three-page book: three
half-planes identified along their common boundary line. Its link is
the graph with two vertices joined by three edges. Consequently its
local second integral homology has rank two. At an interior point of a
two-manifold this local homology has rank one, and at a boundary point
it has rank zero. Thus points of $C$ are not manifold points even when
boundary is permitted. Every point of $X\setminus C$ is a
two-manifold point, possibly with boundary: the square corners are
topological boundary points, and the top circle of the cylinder is a
boundary circle. Hence $C$ is exactly the intrinsic non-manifold locus
of $X$.

Suppose $h:|K|\to X$ were a finite triangulation restricting to the
prescribed identity triangulation on a subcomplex $L$ with $|L|=Y$,
possibly after linear subdivision. In a finite polyhedron the local
germ is constant along the interior of each simplex, by its product
description using the link. Therefore the non-manifold locus is a
union of open simplices. Since $h^{-1}(C)$ is closed, every face of
such a simplex is included too, and it is a finite subpolyhedron.
Its intersection with $|L|$ is then a subpolyhedron of $Y$. Since
$h|_Y$ is the identity, this says that the round circle $C$ is a finite
union of straight simplices. It has no two-dimensional simplex and
cannot contain a nondegenerate straight segment, since a line meets
the circle in at most two points. A finite union of points cannot be
the circle. This is the contradiction.

For the first repair, apply the finite-family definable triangulation
theorem, for example \cite[Theorem 2.2(1), with the constant
function]{Shiota2010}. For the second, let $h:|K|\to X$ be an ambient
triangulation compatible with $Y$, with $Y=h(|L|)$ for a subcomplex
$L$. Let $\tau:|P|\to Y$ be the prescribed triangulation and assume
$h^{-1}\tau:|P|\to|L|$ is PL. Take finite subdivisions on which this
map is simplicial. The resulting subdivision of $L$ extends to $K$:
induct over its remaining simplices, coning each already subdivided
boundary from an interior point. On the subdivided subcomplex the
ambient realization is exactly $\tau$ under this simplicial
identification. This proves the asserted extension with its precise
compatibility hypothesis.
\end{proof}

We checked the sentence directly on the original image of 121/065;
its assertion that every prescribed $\mathcal M$-triangulation of $Y$
is induced is present without a doubt mark. The local homology
certificate in \texttt{code/check\_A307\_links.py} computes the ranks
for one, two and three pages and records the actual output. This
does not replace the topological argument above.

\begin{saresult}{A.308}{Refuted}{The compact closure and gluing axioms do not force tameness}
Let $I=[-1,1]\subset\R$ and let $\mathcal C_n$ be the family of
all closed subsets of $I^n$. Let $\mathcal C$ be the category of
compact spaces admitting an embedding in a finite cube, with all
continuous maps as morphisms. This model satisfies the compact
axioms of 121/054--064, including finite intersections, all the
specified direct and inverse images, the field operations on their
compact domains, interval injectivity, and the zero-set condition.
It also satisfies the compact-neighbourhood and gluing conditions
of 121/070--073, using all compact subsets of $\R^n$ there.
Amalgams along two closed embeddings exist, have the asserted
intersection, and are preserved by the functor to compact spaces.
More generally a pushout along one closed embedding and an
arbitrary morphism exists in this category.

Nevertheless this model contains the Cantor set, is not a system
of compact definable sets with o-minimal open core, and fails
finite triangulability. Thus these axioms, with intersections
and the corrected interval, do not imply the proposed
characterization. The extra triangulation claims of 121/065
must be separated: finite triangulability is independent of
these axioms, while its arbitrary-prescribed-extension clause
is false even for compact semialgebraic sets by A.307.
The exact additional completion criterion is A.306.
\end{saresult}
\begin{proof}
The graph of a continuous map between compact Hausdorff spaces is
closed. Conversely, a map between such spaces with closed graph
is continuous: the projection of its intersection with the inverse
image of a closed subset of the target is compact, hence closed
in the source. Consequently the graph description of morphisms
in 121/056--058 gives precisely the category just stated.
A continuous bijection from a compact space to a Hausdorff space
is a homeomorphism: it maps closed sets to compact, hence closed,
sets. The forgetful functor is therefore faithful and conservative.
It can be made transportable by taking all spaces homeomorphic
to the indicated cube subsets.

Finite products are compact subsets of a finite product cube.
Equalizers are closed subsets of their sources, so all finite
limits exist and are preserved. Continuous images are compact;
thus the image factorization is present. A continuous surjection
of compact Hausdorff spaces is a quotient map, since it is closed.
It follows directly that it is the coequalizer of its kernel pair:
a continuous map constant on its fibres descends continuously
by the quotient topology. The same argument applies after any
base change. Finite unions and intersections of compact subsets
are compact. The coordinate maps in the manuscript, and all
restricted rational scalar, midpoint, and product maps, are
continuous; their graphs, images, and closed-set inverse images
therefore have the required membership. Finite disjoint sums
embed by adding a coordinate taking distinct values. Every object
embeds in a finite power of $I$ by definition, and every point
is allowed.

Here is a full proof of the interval extension property used in
the gluing axiom. Let $A\subset X$ be closed, where $X$ is one
of our compact metrizable spaces, and let $f:A\to[-a,a]$ be
continuous with $a>0$. Set
\[
 A_- =\{x\in A:f(x)\leq-a/3\},\qquad
 A_+ =\{x\in A:f(x)\geq a/3\}.
\]
If both sets are nonempty, they are disjoint closed subsets of
$X$, and
\[
 g(x)=\frac a3\,
 \frac{d(x,A_-)-d(x,A_+)}{d(x,A_-)+d(x,A_+)}
\]
is continuous, has absolute value at most $a/3$, and satisfies
$|f-g|\leq2a/3$ on $A$. The denominator is positive everywhere:
distance zero to a closed set means membership, and the two
sets are disjoint. If only $A_-$ is empty use the constant
$a/3$; if only $A_+$ is empty use $-a/3$; if both are empty
use zero. The same error bound holds in each case.
Apply this construction successively to the residual functions
on $A$, with bounds $a_j=(2/3)^ja$. We obtain continuous
$g_j:X\to[-a_j/3,a_j/3]$. Their series converges uniformly
since $\sum_j a_j/3=a$, so its sum $\bar f$ is continuous,
takes values in $[-a,a]$, and restricts to $f$ because the
residual bound tends to zero. For completeness, continuity of
the uniform limit follows by choosing a tail of sup norm less
than one third of any given error and using continuity of the
finite partial sum for the remaining third. This proves
interval injectivity without an unstated extension theorem.

Every closed subset $A$ of $X$ is a zero set. For nonempty $A$
use $d_A(x)=d(x,A)/(1+d(x,A))$, with values in $I$; for
$A=\varnothing$ use the constant one. The identity zero
function treats $A=X$. Hence the finite zero-set axiom holds
already with a single function.

We verify the pushout explicitly. Write a closed embedding as
$A\subset Y$, let $f:A\to Z$ be continuous, and choose
embeddings $y:Y\hookrightarrow I^p$ and
$z:Z\hookrightarrow I^q$. Extend the coordinates of $z f$
to a continuous map $F:Y\to I^q$ by the preceding construction.
Choose $d:Y\to[0,1]$ with zero set exactly $A$. The two maps
into a finite cube are
\[
 Y\longrightarrow I^{q+1+p},\quad
 u\longmapsto(F(u),d(u),d(u)y(u)),\qquad
 Z\longrightarrow I^{q+1+p},\quad
 v\longmapsto(z(v),0,0).
\]
Their images are compact. Outside $A$ the second coordinate
is positive and the last block, divided by that coordinate,
recovers $y(u)$ and then $u$. On $A$ the first map equals
the second map composed with $f$. Consequently the only
identifications in the union of these images are exactly those
prescribed by the set-theoretic pushout $Y\amalg_A Z$.
The map from the compact space $Y\amalg Z$ to that union is
a continuous surjection onto a Hausdorff space, hence a
quotient map. Thus the union has the topological pushout
property, and lies in $\mathcal C$. If $f$ is an embedding,
both summands embed and their intersection is exactly $A$.
All stated preservation properties follow from this explicit
topological realization; gluing continuous maps on a finite
closed cover follows because inverse images of closed sets
are finite unions of closed sets. Pulling back such a cover
gives another finite closed cover, so its gluing property is
stable under arbitrary base change.

For 121/070--073 take $C_n$ to be every compact subset of
$\R^n$. Finite-coordinate direct images, coordinate
repetitions, products and intersections preserve compactness.
The graph argument gives again precisely continuous maps.
Compact neighbourhoods form a neighbourhood basis on each
object: intersect it with a sufficiently small closed Euclidean
ball. The interiors of closed boxes form a basis in $\R^n$.
The same extension and pushout arguments prove the remaining
claims of that draft. No assertion about quotients by arbitrary
closed equivalence relations, or all finite colimits, is needed.

Finally the middle-thirds Cantor set $K\subset[0,1]$ is a
closed subset and hence an object. It is infinite and totally
disconnected: two distinct points are separated by one of the
deleted middle-third intervals, whose complementary parts
give a separation of $K$. It cannot be a finite polyhedron.
A finite polyhedron with a positive-dimensional simplex
contains a connected nondegenerate segment, while a finite
polyhedron having only zero-dimensional simplices is finite.
Also $K$ is not a finite union of points and intervals. Were
these compact sets the compact definable sets of an expansion
with o-minimal open core, its definable open set $\R\setminus K$
would belong to that open core, and its complement $K$ would
be definable there. This contradicts o-minimality. This proves
the counterexample and all its required hypotheses.
\end{proof}
